\begin{RequAnswer}{6.1}
$f(n)=O(g(n))$ means that $f(n)$ grows no faster than $g(n)$.
Another way of phrasing it is saying that $g(n)$ is an upper bound on $f(n)$.
{\em If you said that it means $g(n)$ grows faster than $f(n)$, you are not quite correct.
Go reread the definition and make sure you are clear on this point!}
$f(n)=\Theta(g(n))$ means that $f(n)$ and $g(n)$ grow at the same rate.
Another way of phrasing it is saying that $g(n)$ is a tight bound on $f(n)$.
\end{RequAnswer}
\begin{RequAnswer}{6.2}
(a) No. It means there is some $n_0$ such that whenever $n\geq n_0$, $f(n)\leq c g(n)$ (where $c>0$ is some constant).
So there are two problems with saying it means $f(n)\leq g(n)$. First, there is a constant involved.
So $f(n)\leq c\, g(n)$, not just $f(n)\leq g(n)$. Second, it does not have to hold for all values of $n$,
but only for (i.e. $n\geq n_0$).
(b) No. Repeating what was said in the previous part, it has to hold for all values of $n$ at least as large as some
given value $n_0$, but it does not have to hold for smaller values of $n$.
(c) That is certainly possible. For instance, if $f(n)=100n$ and $g(n)=n^2$, $g(10)=100<1000=f(10)$,
but hopefully it is clear that $f(n)=O(g(n))$.
\end{RequAnswer}
\begin{RequAnswer}{6.3}
No. For instance, if $f(n)=23n$ and $g(n)=n^2$,  then $f(n)=O(g(n))$, but  $f(n)\not=\Theta(g(n))$.
\end{RequAnswer}
\begin{RequAnswer}{6.4}
Yes. $f(n)=\Theta(g(n))$ if and only if $f(n)=O(g(n))$ and $f(n)=\Omega(g(n))$.
So if $f(n)=\Theta(g(n))$, then $f(n)=O(g(n))$.
\end{RequAnswer}
\begin{RequAnswer}{6.5}
If $f(n)=o(g(n))$, then we know that $f(n)$ grows {\em slower than} $g(n)$.
If $f(n)=O(g(n))$, then $f(n)$ grows {\em no faster than} $g(n)$.
That is, $f(n)$ either grows slower than $g(n)$ or they grow at the same rate.
\end{RequAnswer}
\begin{RequAnswer}{6.6}
More. If you know that $f(n)=\Theta(g(n))$, then you know that they grow at the same rate.
But if you only know that $f(n)=O(g(n))$, it is possible that the grow at the same rate, but it is also possible
that $f(n)$ grows slower than $g(n)$.
\end{RequAnswer}
\begin{RequAnswer}{6.7}
Proof 1:
Notice that if $n\geq 1$, $7n^3+4n^2-8n+27\leq 7n^3+4n^2+27 \leq 7n^3+4n^3+27n^3=38n^3$.
By definition of Big-O, $7n^3+4n^2-8n+27=O(n^3)$.

\noindent
Proof 2: Notice that
$\displaystyle \lim_{n\rightarrow\infty} \frac{7n^3+4n^2-8n+27}{n^3} =
\lim_{n\rightarrow\infty} \frac{7n^3}{n^3}+\frac{4n^2}{n^3}-\frac{8n}{n^3}+\frac{27}{n^3}
=\lim_{n\rightarrow\infty} 7+\frac{4}{n}-\frac{8}{n^2}+\frac{27}{n^3}
=7+0+0+0=7.$
Thus, $7n^3+4n^2-8n+27=\Theta(n^3)$ by Theorem~\ref{thm:limits} (part 3).
By Theorem~\ref{thm:theta}, $7n^3+4n^2-8n+27=O(n^3)$.
\end{RequAnswer}
\begin{RequAnswer}{6.8}
Notice that
$\displaystyle \lim_{n\rightarrow\infty} \frac{3^n}{3.1^n} = \lim_{n\rightarrow\infty} \left(\frac{3}{3.1}\right)^n=0$
since $\frac{3}{3.1}<1$. By Theorem~\ref{thm:limits} (part 1), $3^n=o(3.1^n)$.
\end{RequAnswer}
\begin{RequAnswer}{6.9}
Notice that both functions have a factor of $n$. Since $\log n$ grows faster than $c$ (which doesn't grow at all since
it is a constant), $n\log n$ grows faster than $c n$.

It is actually easy to prove this formally:
$\displaystyle\lim_{n\rightarrow\infty} \frac{n\log n}{c\, n} = \lim_{n\rightarrow\infty} \frac{\log n}{c} =\infty$,
so by Theorem~\ref{thm:limits}, $n\log n=\omega(c n)$. In other words, $n\log n$ grows faster than $c n$.

\end{RequAnswer}
\begin{RequAnswer}{6.10}
No! The function may grow slowly, but {\em it is still growing}.
So you are multiplying a function by another function that is growing (even if slowly), which grows faster than 1.
Thus, $f(n)\log n$ definitely grows faster than $f(n)$.

As with the previous question, a simple limit computation gives a clear proof of this fact.
$\displaystyle\lim_{n\rightarrow\infty} \frac{f(n)\log n}{f(n)} = \lim_{n\rightarrow\infty} \frac{\log n}{1} =\infty$,
so by Theorem~\ref{thm:limits}, $f(n)\log n=\omega(f(n))$. In other words, $f(n)\log n$ grows faster than $f(n)$.

\end{RequAnswer}
\begin{RequAnswer}{6.11}
$8675309$,
$7 \log_{10} n$ $\sim$ $\log_3 n$,
$27 n$,
$7n\log_2 n$,
$n^2+n+1$ $\sim$ $3n^2$,
$n^3$ $\sim$ $n^3+n^2\log_{e} n$,
$2^n$,
$7^n$,
$n!$,
$n^n$
\end{RequAnswer}
\begin{RequAnswer}{6.12}
It is unclear what types of machines the algorithms were run on. If one was run on a cellphone and another
on a supercomputer, it is definitely not a fair comparison.
If they were run on the same machine and on the same data, then we can compare them.
\end{RequAnswer}
\begin{RequAnswer}{6.13}
No. If one always takes 3 times as long, it's just a constant multiple difference, so they have
the same growth rate. E.g. if one takes $f(n)$ time, the other takes $3f(n)$ time, and these
have the same growth rate.
Remember, when we are talking about growth rates, the constant multiples do not matter.
\end{RequAnswer}
\begin{RequAnswer}{6.14}
Since $g(n)$ grows faster than $f(n)$, as $n$ increases, algorithm $B$ will take (relatively) longer than $A$.
In other words, algorithm $A$ is faster.
{\em Remember: Faster growth rate of computational complexity means slower algorithm!}
\end{RequAnswer}
\begin{RequAnswer}{6.15}
You cannot even read all of the inputs in $\Theta(\log n)$ time, so this is clearly ludicrous.
It would have to take {\em at least} $\Omega(n)$ time
(and even that is not attainable, but we aren't ready for that proof yet).
\end{RequAnswer}
\begin{RequAnswer}{6.16}
This would be an awful algorithm since searching should not take more than $\Theta(n)$ time.
\end{RequAnswer}
\begin{RequAnswer}{6.17}
Better because \verb+insertionSort+ takes $\Theta(n^2)$ time and $n^{1.5}=o(n^2)$.
\end{RequAnswer}
\begin{RequAnswer}{6.18}
One (or both) of the loops could only execute a constant number of times.
Also, maybe the index variable(s) aren't just incremented/decremented by 1 each time through the loop.
If instead the variables are doubled or halved, for instance, the complexity might involve a logarithm or something
(e.g. maybe $\Theta(n\log n)$).
Here's a third (for free!): Maybe the outer loop executes $n$ times and the inner loop executes $m$ times.
Then the complexity is $\Theta(n\, m)$.
\end{RequAnswer}
\begin{RequAnswer}{6.19}
(a) Generally speaking, $B$ is faster. Since we are given $\Theta$ bounds on the complexities, we know
the exact growth rates and $n\log n = o(n^2)$, so $B$ is definitely a faster algorithm.
We know that $B$ is faster for large enough input. We do not know if is is faster for small inputs.
(b) As mentioned in the previous part, for small values of $n$, $A$ could possibly be faster because the
constants involved in $B$ might be much larger.
\end{RequAnswer}
\begin{RequAnswer}{6.20}
This is a trick question! The answer is unclear. We know that $A$ has complexity $\Theta(n\log n)$, but
we are only given a Big-O bound on the complexity of $B$. It is possible that $B$ has complexity
$\Theta(n)$ and it is faster, or that it has complexity $\Theta(n\log n)$ and it is the same as $A$,
or that it has complexity $\Theta(n^2)$ and it is slower than $A$.
\end{RequAnswer}
\begin{RequAnswer}{6.21}
Memory usage;
how complicated the algorithms are to implement, maintain, and debug;
the constants; how convinced you are of the correctness of each algorithm.
\end{RequAnswer}
